\section{Conclusion}
\label{sec:conclusion}

We have presented the Squiggly Line Conjecture: randomly initialized complex functions tend to be geometrically closer to random data points than simpler functions through increased spatial coverage. This offers a geometric perspective on overparameterized model behavior.

We observe that the $w^{-\alpha}$ scaling law appears across polynomials, neural networks, and quantum systems, with approximately dimension-independent behavior. Initial geometric advantage may contribute to faster convergence, potentially offering unified explanations for phenomena like the lottery ticket hypothesis and overparameterization benefits.

\textbf{Limitations and future work}: While effect sizes are typically 10-50\% improvements and benefits saturate around $w = \sqrt{n}$, future directions include architectural optimization for coverage, quantum machine learning applications, and biological validation.

\section*{Acknowledgments}

We thank the machine learning community for valuable discussions and insights that shaped this work.

\section*{Broader Impact}

This theoretical work may advance understanding of neural network initialization and overparameterization through geometric principles, with potential applications to more efficient architectures and training procedures.
